\documentclass[10pt]{article}
\usepackage{cvstyle}

\begin{document}

\cvheader{Postdoctoral Researcher · Computational Science Lab, University of Amsterdam}{%
  {\color{cvteal}\faEnvelope}~\href{mailto:c.j.hourican@uva.nl}{c.j.hourican@uva.nl}\quad
  {\color{cvteal}\faGlobe}~\href{https://cillianhourican.com}{cillianhourican.com}\quad
  {\color{cvteal}\faLinkedin}~\href{https://www.linkedin.com/in/cillian-hourican-5460559a}{in/cillian-hourican}\quad
  {\color{cvteal}\faGithub}~\href{https://github.com/CillianHourican}{CillianHourican}\quad
  {\color{cvteal}\aiOrcid}~\href{https://orcid.org/0009-0002-3161-0736}{0009-0002-3161-0736}%
}

\vspace{6pt}
{\RaggedRight
Mathematician and computational scientist working across statistics, machine learning, networks and information theory. I build rigorous models of complex, interconnected systems --- higher-order interactions and causality in health and ageing --- and explain them clearly. Currently a postdoctoral researcher in the Computational Science Lab at the University of Amsterdam; PhD defense scheduled for 27 November 2026.\par}

\section{Education}
\entry{2022 -- present}{%
  \etitle{PhD in Computational Science}\\[1pt]
  \eplace{University of Amsterdam}\\[2pt]
  \edetail{Dissertation: \emph{Dynamic Symptom Networks to Study Multimorbidity} --- at the intersection of information theory, network science and causality. Supervised by Dr.\ Rick Quax. \textbf{Defense scheduled for 27 November 2026.}}}
\entry{2021}{%
  \etitle{MSc Computational Science}\ \edetail{(GPA 8/10)}\\[1pt]
  \eplace{University of Amsterdam}\\[2pt]
  \edetail{Thesis: non-intrusive uncertainty quantification of a 3D in-stent restenosis model using active learning and adaptive surrogate modelling.}}
\entry{2018 -- 2019}{%
  \etitle{Research Student, Optimisation \& Machine Learning}\\[1pt]
  \eplace{Trinity College Dublin}\\[2pt]
  \edetail{Online convex optimisation and bandits; deep-learning coursework (CNNs, image classification).}}
\entry{2018}{%
  \etitle{BA (Mod.) Mathematics}\ \edetail{First-Class Honours}\\[1pt]
  \eplace{Trinity College Dublin}\\[2pt]
  \edetail{Final-year project: estimating social networks with stochastic block models and multiplex networks.}}

\section{Publications}
\begin{pubs}
  \item \me{Hourican, C.}, Peeters, G., Melis, R.J.F., Kok, A., van Schoor, N.M., Wezeman, S., Lees, M., Olde Rikkert, M.G.M., \& Quax, R. (2026). \emph{Interpreting higher-order dependence in multimorbidity using cohort data: a partial information decomposition approach.} \venue{PLOS Computational Biology}. \doilink{10.1371/journal.pcbi.1014386}
  \item Peeters, G.\textsuperscript{*}, \me{Hourican, C.}\textsuperscript{*}, Lees, M., Quax, R., Melis, R., Kok, A., van Schoor, N.M., \& Olde Rikkert, M. (2026). \emph{Symptoms provide predictive information on daily functioning beyond signs and diseases in middle-aged and older adults, particularly those with multimorbidity.} \venue{Age and Ageing}. \doilink{10.1093/ageing/afag046}\ {\small\color{cvmuted}(*joint first authors)}
  \item Li, J., Bosch, J.A., Rydin, A.O., \me{Hourican, C.}, Koloi, A., Tassi, S.C., Mishra, P.P., Mishra, B.H., K\"ah\"onen, M., Lehtim\"aki, T., Raitakari, O.T., Laaksonen, R., Keltikangas-J\"arvinen, L., Juonala, M., \& Quax, R. (2025). \emph{A multilayer network analysis of cardiovascular--depression comorbidity reveals symptom-specific molecular biomarkers.} \venue{Psychological Medicine}. \doilink{10.1017/S0033291725102109}
  \item \me{Hourican, C.}, Li, J., Mishra, P.P., Lehtim\"aki, T., Mishra, B.H., K\"ah\"onen, M., Raitakari, O.T., Laaksonen, R., Keltikangas-J\"arvinen, L., Juonala, M., \& Quax, R. (2024). \emph{Efficient search algorithms for identifying synergistic associations in high-dimensional datasets.} \venue{Entropy}. \doilink{10.3390/e26110968}
  \item Gehlen, J., Li, J., \me{Hourican, C.}, Tassi, S.C., Mishra, P.P., Lehtim\"aki, T., K\"ah\"onen, M., Raitakari, O., Bosch, J.A., \& Quax, R. (2024). \emph{Bias in O-information estimation.} \venue{Entropy}. \doilink{10.3390/e26100837}
  \item Koloi, A., Loukas, V.S., \me{Hourican, C.}, Sakellarios, A.I., Quax, R., Mishra, P.P., Lehtim\"aki, T., Raitakari, O.T., Papaloukas, C., Bosch, J.A., M\"arz, W., \& Fotiadis, D.I. (2024). \emph{Predicting early-stage coronary artery disease using machine learning and routine clinical biomarkers improved by augmented virtual data.} \venue{European Heart Journal --- Digital Health}. \doilink{10.1093/ehjdh/ztae049}
  \item \me{Hourican, C.}, Peeters, G., Melis, R.J.F., Wezeman, S.L., Gill, T.M., Olde Rikkert, M.G.M., \& Quax, R. (2023). \emph{Understanding multimorbidity requires sign--disease networks and higher-order interactions, a perspective.} \venue{Frontiers in Systems Biology}. \doilink{10.3389/fsysb.2023.1155599}
\end{pubs}

\section{Research \& Professional Experience}
\entry{Apr 2026 -- present}{%
  \etitle{Postdoctoral Researcher}\\[1pt]
  \eplace{Computational Science Lab, University of Amsterdam}\\[2pt]
  \edetail{Higher-order interactions, networks and causality in health and complex systems.}}
\entry{2017}{%
  \etitle{Summer Research Student}\\[1pt]
  \eplace{Trinity College Dublin}\\[2pt]
  \edetail{Changepoint-detection algorithms: frequentist and Bayesian (Gibbs sampling) approaches.}}
\entry{2016}{%
  \etitle{Summer Research Student}\\[1pt]
  \eplace{LM Ericsson, Ireland}\\[2pt]
  \edetail{Graph-theoretic methods for self-organising networks.}}

\section{Teaching}
\entry{2023 -- 2025}{\etitle{Introduction to Computational Science}\ \edetail{— MSc, teaching assistant}\\ \eplace{University of Amsterdam}}
\entry{2022 / 2023}{\etitle{Scientific Data Analysis}\ \edetail{— BSc Informatica, teaching assistant}\\ \eplace{University of Amsterdam}}
\entry{2021 -- 2023}{\etitle{Seminars Computational Science}\ \edetail{— MSc, co-organiser \& teaching assistant}\\ \eplace{University of Amsterdam}}
\entry{2019 -- 2021}{\etitle{Modelling System Dynamics}\ \edetail{— MSc, teaching assistant (Vensim, AnyLogic)}\\ \eplace{University of Amsterdam}}
\entry{2017 -- 2019}{\etitle{Statistics \& Probability tutorials}\ \edetail{— BSc, teaching assistant (R)}\\ \eplace{Trinity College Dublin}}

\section{Student Supervision}
\subsection{Daily supervisor}
\begin{itemize}
  \item \textbf{Jacco van Wijk} (MSc, 2023) --- \emph{Predicting the Accuracy of Dynamic Bayesian Network Structure Learning Algorithms}.
  \item \textbf{K.\ Ullah} (MSc, 2026) --- \emph{Survival of the Fittest Structure: Evolutionary Learning of Synergistic Dependencies}.
  \item \textbf{N.\ Kashyap} (MSc Information Studies, 2024) --- \emph{Predictive Modeling of Symptom Reports and Changes in Functional Outcomes in Elderly Suffering from Multimorbidity}.
  \item \textbf{Jay Sitabi} (BSc Informatica).
  \item \textbf{Jasper van Willigen} (BSc Informatica).
\end{itemize}
\subsection{Second supervisor / assessor}
\begin{itemize}
  \item \textbf{Johanna Gehlen} (MSc, 2023) --- \emph{Bias in O-Information Estimation} (led to a peer-reviewed publication in \venue{Entropy}).
  \item \textbf{Noah van de Bunt} (MSc, 2023) --- \emph{A Bayesian Network Study to Uncover Bridging Mechanisms Causing the Comorbidity of Major Depressive Disorder and Cardio-Metabolic Diseases}.
  \item \textbf{D.\ Klein Leunk} (MSc Computational Science).
  \item \textbf{Leon Kuiper} (MSc Computational Science).
  \item \textbf{Sofia T\'et\'e Garcia Ramos Nunes} (MSc Computational Science).
  \item \textbf{Seda den Boer} (MSc Computational Science).
\end{itemize}

\section{Talks \& Conferences}
\begin{itemize}
  \item \textbf{Oral} --- ``Removing unnecessary hyperedges in information-theoretic hypergraphs through O-information'', NetSci, Vienna, 2023.
  \item \textbf{Poster} --- International Conference on Network Science (NetSci), 2025.
  \item \textbf{Poster} --- Dutch Computational Science Day (DUCOMS), Utrecht, 2023.
  \item \textbf{Poster} --- Dutch NetSci, Delft, 2023.
  \item \textbf{Participant} --- Gateway to Global Aging research hackathon, USC, 2023.
  \item \textbf{Attendee} --- Decomposing Multivariate Information in Complex Systems (demics23), Dresden, 2023.
\end{itemize}

\section{Technical Skills}
\skillrow{Languages}{Python, R}
\skillrow{Methods}{Machine learning \& deep learning · statistical \& predictive modelling · uncertainty quantification · surrogate modelling}
\skillrow{Fields \& tools}{Network science · information theory · causal inference · complex systems · Vensim · AnyLogic}

\end{document}
